\documentclass[10pt,a4paper]{amsart}
\usepackage{amsmath,amssymb,amsthm,graphicx}
% Retained minimal preamble for the Stacks Project chapter more-algebra.tex (More on Algebra).
% The chapter file pinned-original-source.tex carries only \input{preamble} on its first line, so no
% preamble text lives inside the chapter itself.  The Stacks Project uses one single shared preamble for
% every chapter; the snapshot stages it identically as sources/stacks-chow/preamble.tex and
% sources/stacks-etale-cohomology/preamble.tex.  The block below is copied mechanically, line for line,
% from sources/stacks-chow/preamble.tex lines 14-18 (the Xy-pic setup) and 154-217 (the native theorem
% declarations together with the mathematical macros).  Nothing was added, removed or redefined; the
% wrapper class and the amsmath/amssymb/amsthm/graphicx base packages are supplied by the assembler.
% Chapter: more-algebra.tex -- More on Algebra.
% For commutative diagrams we use Xy-pic
\usepackage[all]{xy}

% We use 2cell for 2-commutative diagrams.
\xyoption{2cell}

% Theorem environments.
%
\theoremstyle{plain}
\newtheorem{theorem}[subsection]{Theorem}
\newtheorem{proposition}[subsection]{Proposition}
\newtheorem{lemma}[subsection]{Lemma}

\theoremstyle{definition}
\newtheorem{definition}[subsection]{Definition}
\newtheorem{example}[subsection]{Example}
\newtheorem{exercise}[subsection]{Exercise}
\newtheorem{situation}[subsection]{Situation}

\theoremstyle{remark}
\newtheorem{remark}[subsection]{Remark}
\newtheorem{remarks}[subsection]{Remarks}

\numberwithin{equation}{subsection}

% Macros
%
\def\lim{\mathop{\mathrm{lim}}\nolimits}
\def\colim{\mathop{\mathrm{colim}}\nolimits}
\def\Spec{\mathop{\mathrm{Spec}}}
\def\Hom{\mathop{\mathrm{Hom}}\nolimits}
\def\Ext{\mathop{\mathrm{Ext}}\nolimits}
\def\SheafHom{\mathop{\mathcal{H}\!\mathit{om}}\nolimits}
\def\SheafExt{\mathop{\mathcal{E}\!\mathit{xt}}\nolimits}
\def\Sch{\mathit{Sch}}
\def\Mor{\mathop{\mathrm{Mor}}\nolimits}
\def\Ob{\mathop{\mathrm{Ob}}\nolimits}
\def\Sh{\mathop{\mathit{Sh}}\nolimits}
\def\NL{\mathop{N\!L}\nolimits}
\def\CH{\mathop{\mathrm{CH}}\nolimits}
\def\proetale{{pro\text{-}\acute{e}tale}}
\def\etale{{\acute{e}tale}}
\def\QCoh{\mathit{QCoh}}
\def\Ker{\mathop{\mathrm{Ker}}}
\def\Im{\mathop{\mathrm{Im}}}
\def\Coker{\mathop{\mathrm{Coker}}}
\def\Coim{\mathop{\mathrm{Coim}}}

% Boxtimes
%
\DeclareMathSymbol{\boxtimes}{\mathbin}{AMSa}{"02}

%
% Macros for moduli stacks/spaces
%
\def\QCohstack{\mathcal{QC}\!\mathit{oh}}
\def\Cohstack{\mathcal{C}\!\mathit{oh}}
\def\Spacesstack{\mathcal{S}\!\mathit{paces}}
\def\Quotfunctor{\mathrm{Quot}}
\def\Hilbfunctor{\mathrm{Hilb}}
\def\Curvesstack{\mathcal{C}\!\mathit{urves}}
\def\Polarizedstack{\mathcal{P}\!\mathit{olarized}}
\def\Complexesstack{\mathcal{C}\!\mathit{omplexes}}
% \Pic is the operator that assigns to X its picard group, usage \Pic(X)
% \Picardstack_{X/B} denotes the Picard stack of X over B
% \Picardfunctor_{X/B} denotes the Picard functor of X over B
\def\Pic{\mathop{\mathrm{Pic}}\nolimits}
\def\Picardstack{\mathcal{P}\!\mathit{ic}}
\def\Picardfunctor{\mathrm{Pic}}
\def\Deformationcategory{\mathcal{D}\!\mathit{ef}}
\usepackage{cleveref}
\crefname{originalremark}{original remark}{original remarks}
\Crefname{originalremark}{Original remark}{Original remarks}
\makeatletter
\newlabel{example-cohomology-complex-tensored}{{061Z}{source}{}{sourcepointer-1}{}}
\newlabel{example-cohomology-complex-tensored@cref}{{[originalremark][0][]061Z}{[1][0][]source}{}{}{}}
\makeatother
\crefname{originallemma}{original lemma}{original lemmas}
\Crefname{originallemma}{Original lemma}{Original lemmas}
\makeatletter
\newlabel{lemma-check-perfect-pointwise}{{068V}{source}{}{sourcepointer-2}{}}
\newlabel{lemma-check-perfect-pointwise@cref}{{[originallemma][0][]068V}{[1][0][]source}{}{}{}}
\makeatother
\crefname{originalproposition}{original proposition}{original propositions}
\Crefname{originalproposition}{Original proposition}{Original propositions}
\makeatletter
\newlabel{algebra-proposition-regular-finite-gl-dim}{{00O7}{source}{}{sourcepointer-3}{}}
\newlabel{algebra-proposition-regular-finite-gl-dim@cref}{{[originalproposition][0][]00O7}{[1][0][]source}{}{}{}}
\makeatother
\title{Perfectness of a pseudo-coherent complex over a flat local extension with regular fibre, from one fibre computation at the closed point}
\author{The Stacks Project Authors}
\date{stacks-project:more-algebra.tex:local-snapshot-20261012; GNU FDL 1.2 or later}
\begin{document}
\maketitle
\setcounter{subsection}{1}
\noindent\textit{Editorial note.} All delivered mathematical text is the author's own, copied from the pinned Stacks Project chapter file more-algebra.tex; no statement or proof step was rewritten, reordered or removed. The counted statement is the labeled lemma environment at source lines 20228-20244 (canonical tag 09PC) and the counted proof is the proof environment at source lines 20246-20273, which immediately follows it and is its own proof.\par\medskip
\noindent\textit{Editorial note.} The proof refers three times: to example-cohomology-complex-tensored and lemma-check-perfect-pointwise inside this chapter, and to algebra-proposition-regular-finite-gl-dim of another chapter. All three are navigation pointers with their roles recorded in the external reference reasons; the displayed computation of the fibre at the closed point and the degree estimate are inside the counted span.\par\medskip
\noindent\textit{Editorial note.} The sentence at source line 20226, announcing that the lemma can be used to find perfect complexes over regular local rings, is outside the counted spans.\par\medskip
\noindent\textit{Editorial note.} The chapter file pinned-original-source.tex carries only an input of the shared preamble on its first line, so no preamble text lives inside the chapter. The retained file sources/stacks-more-algebra/retained-preamble.tex is a mechanical, line-for-line copy of the project's single shared preamble as staged in the snapshot (Xy-pic setup plus the native theorem declarations and mathematical macros). Nothing was added, removed or redefined, and the excerpt serves only as preamble, never as a span coordinate.\par\medskip
\noindent\textit{Editorial note.} The source URL is the source-level Stacks Project URL, as required when more than one card is produced from one source; the per-claim tag is carried by claim\_key and by this note. The counted statement and proof cite no bibliography entry, so this card needs no bibliography.\par\medskip
\noindent\textit{Allowed standard prerequisites (editorial).}\begin{itemize}
\item The derived tensor product with the residue field at the closed point of a local ring detects perfectness of a pseudo-coherent complex, through the vanishing of the cohomology of that fibre outside the tor amplitude range. invoked at proof line 20272, where the vanishing obtained for j outside [a - d, b] is turned into the conclusion. Converts the single fibre computation at the closed point into perfectness over the local ring B.
\item A regular local ring has finite global dimension equal to its dimension, so the fibre B over m\_A B has global dimension d. invoked at proof lines 20267-20269 by reference to the algebraic result that regular rings have finite global dimension. Supplies the bound d that shifts the tor amplitude range and makes the spectral sequence argument terminate.
\item For a bounded above complex of flat modules the ordinary tensor product computes the derived tensor product, and the spectral sequence of the example on cohomology of a tensor product of complexes relates the cohomology of the fibre to that of K tensor\_A kappa(m\_A). used at proof lines 20257-20261 and 20265-20267, where the three equalities of the display are justified and the spectral sequence is applied. Produces the identification K tensor\textasciicircum{}L\_B kappa(m\_B) = (K tensor\_A kappa(m\_A)) tensor\textasciicircum{}L over B over m\_A B kappa(m\_B) and hence the degree estimate.
\end{itemize}\medskip
\ifdefined\hypertarget\hypertarget{sourcepointer-1}{}\fi
\noindent\textit{Source reference.} Original remark 061Z: sources/stacks-more-algebra/pinned-original-source.tex, source lines 15159--15159. The reference names the spectral sequence of a tensor product of complexes used at proof lines 20265-20267 to pass from the cohomology bound on the fibre over A to the bound on the fibre over B; the identification of the fibre and the degree bookkeeping around it are inside the counted span.\par\medskip
\ifdefined\hypertarget\hypertarget{sourcepointer-2}{}\fi
\noindent\textit{Source reference.} Original lemma 068V: sources/stacks-more-algebra/pinned-original-source.tex, source lines 20032--20032. The reference is the criterion invoked at proof line 20272 to conclude perfectness from the vanishing of the cohomology of the fibre; the vanishing itself is proved inside the counted span.\par\medskip
\ifdefined\hypertarget\hypertarget{sourcepointer-3}{}\fi
\noindent\textit{Source reference.} Original proposition 00O7: sources/stacks-chow/tags.txt, source lines 885--885. The reference supplies the standard fact that a regular ring has finite global dimension, used at proof lines 20267-20269 to bound the global dimension of the regular fibre by its dimension d; the spectral sequence argument that consumes the bound is inside the counted span.\par\medskip
\begin{lemma}
\label{lemma-perfect-over-regular-local-ring}
Let $A \to B$ be a local ring homomorphism.
Let $a, b \in \mathbf{Z}$. Let $d \geq 0$.
Let $K^\bullet$ be a complex of $B$-modules. Assume
\begin{enumerate}
\item the ring map $A \to B$ is flat,
\item the ring $B/\mathfrak m_AB$ is regular of dimension $d$,
\item $K^\bullet$ is pseudo-coherent as a complex of $B$-modules, and
\item $K^\bullet$ has tor amplitude in $[a, b]$ as a complex
of $A$-modules, in fact it suffices if
$H^i(K^\bullet \otimes_A^\mathbf{L} \kappa(\mathfrak m_A))$
is nonzero only for $i \in [a, b]$.
\end{enumerate}
Then $K^\bullet$ is perfect as a complex of $B$-modules
with tor amplitude in $[a - d, b]$.
\end{lemma}

\begin{proof}
By (3) we may assume that $K^\bullet$ is a bounded above complex of finite free
$B$-modules. We compute
\begin{align*}
K^\bullet \otimes_B^{\mathbf{L}} \kappa(\mathfrak m_B)
& = K^\bullet \otimes_B \kappa(\mathfrak m_B) \\
& = (K^\bullet \otimes_A \kappa(\mathfrak m_A))
\otimes_{B/\mathfrak m_A B} \kappa(\mathfrak m_B) \\
& = (K^\bullet \otimes_A \kappa(\mathfrak m_A))
\otimes^{\mathbf{L}}_{B/\mathfrak m_A B} \kappa(\mathfrak m_B)
\end{align*}
The first equality because $K^\bullet$ is a bounded above complex
of flat $B$-modules. The second equality follows from basic 
properties of the tensor product. The third equality holds because
$K^\bullet \otimes_A \kappa(\mathfrak m_A) =
K^\bullet/ \mathfrak m_A K^\bullet$ is a bounded above complex
of flat $B/\mathfrak m_A B$-modules. Since $K^\bullet$ is a bounded
above complex of flat $A$-modules by (1), the cohomology modules $H^i$
of the complex $K^\bullet \otimes_A \kappa(\mathfrak m_A)$ are nonzero only
for $i \in [a, b]$ by assumption (4). Thus the spectral sequence
of Example \ref{example-cohomology-complex-tensored} and the
fact that $B/\mathfrak m_AB$ has finite global dimension $d$
(by (2) and
Algebra, Proposition \ref{algebra-proposition-regular-finite-gl-dim})
shows that $H^j(K^\bullet \otimes_B^{\mathbf{L}} \kappa(\mathfrak m_B))$
is zero for $j \not \in [a - d, b]$.
This finishes the proof by Lemma \ref{lemma-check-perfect-pointwise}.
\end{proof}

\end{document}
